\section{Stelle di neutroni: Intro storica}

\begin{frame}{Brief History of NS}
    \begin{itemize}
        \item 1931 - Landau article on maximus mass of compact star: idea of compact star as giant nucleus
        \item 1932 - Chadwick discovers neutron
        \item 1934 - Article by Baade and Zwicky in which appeared for first time the Neutron stars: speculated aboud collapse of SN to neutron star
        \item 1934 - Oppenheimer, Volkoff studied neutron stars within General relativityi
        \item 1934 - Tolman ''Relativity, Thermodynamics and Cosmology'': fluid sphere in general relativity
        \item 1939 - TOV solution: OV numerically, T making assumption about metric. Maximum mass of star composed on gas of neutrons $0.7\msun{}$, using Tolman analytical solution shown than maxmass con be incresed not further of $1\msun$ including repulsion between neutrons.
        \item 1942 - Duyvendak: Chinise/Japonese records of ''guess star'' close to $\zeta$ Tauris in 1054 AD.
        \item 1942 - Mayall- Oort: Identified 1054 event with SN explosion and Crab nebula in Tauris constellation.
        \item 1958 - Harrison, Wheeler: Unified EOS connection dense matter from WD to NS. They used Empirical mass formula for nmuclei, the liquid drop model, and calculated EOS for mixture of nuclei, electrons and a free gas of neutron minimizing total energy in $\beta$-equilibrium. (Harrison et al 1965).
        \item 1967 - Cocke - Pulsar discovery. Association of star in middle of Crab nebula to Neutron star created by SN explosion.
        \item Cameron 1959 - Study Neutron Stars taking into account interactions between neutrons. Model of skyrme adapted for nuclear matter (Skyrme 1959). Maximum mass $2.0\msun$, radius \SI{8.1}{\kilo\meter}; predicted hyperons at higher densities.
        \item Salpeter 1960 - EOS including hyperons
        \item Ambartsumyan, Saakyan 1960 - Investigates appearence of hyperons in dense neutron star matter for Fermi gas: they found \PSigmaminus appears at density \SI{0.64}{\per\cubic\femto\meter} and \PLambda at \SI{1.27}{\per\cubic\femto\meter} which are 4 times and 8 times the saturation density $n_0=\SI{0.15}{\per\cubic\femto\meter}$
        \item Ambartsumyan, Saakyan 1962 - Mass-Radius relation of NS including hyperon for Fermi gas. EOS for nucleons and hyperons with interaction term proportional to number density squared - hyperons appear in core of massive NS; maximum mass in free case is $0.63\msun$ and radius \SI{11}{\kilo\meter}, in interacting case $1.0\msun$ and radius \SI{5.4}{\kilo\meter}.
\end{itemize}
\end{frame}

\section{Neutron Star Crust}

\subsection{Outer Crust}

\begin{frame}{Atmosphere}
    \begin{itemize}
        \item Transition atmosphere-solid crust: Fermi energy uguale a binding energy of last electron: $E_F=\frac{k_F^2}{2m_e}\approx E_b(Z)=\SI{13.6}{\ev}Z^2$
        \item Spessore atmosfera: da equilibrio idrostatico $\TDy{r}{P}\approx-\frac{P_c}{\Delta R}\approx-\frac{GM\rho_c}{R^2}$, quindi $\frac{\Delta R}{R}\approx \frac{R}{GM}\frac{2E_F}{5m_u}\frac{A}{Z}$; nel caso di stelle di neutroni lo spessore dell'atmosfera sar\'a $\frac{\Delta R}{R}\approx\SI{2.92e-8}(\frac{0.2}{C})ZA$: \SI{0.3}{\milli\meter} per atmosfera di H, \SI{10}{\cm} per atmosfera di Fe.
        \item Determine spectral properties: photon mean free path is smaller than atmospheric thickness.
        \item Atmospheric composition depends upon environment and formation conditions. For binary NS atmosphere will be H accreted from companion; for isolated NS moving in ISM composition will be mainly H; newly born NS will have composition related to material falling back from SN explosion dominated by heavy elements such Silicon.
    \end{itemize}
\end{frame}

\begin{frame}{BPS model} 
    \begin{itemize}
        \item Modello di Baym, Pethick, Sutherland (1971): reticolo di nuclei neutron-rich e gas di elettroni degeneri, assenza di neutroni liberi; $T=0$ equilibrio termodinamico/$\beta$-equilibrio; unica specie nucleare $(A,Z)$ per ogni strato; nuclei in reticolo coulombiano bcc; elettroni degeneri trattati come gas relativistico ideale; neutralit\'a elettrica $n_e=Zn_N$ dove $n_N$ \'e densit\'a dei nuclei; densit\'a barionica $n_b=An_N$; assenza neutroni liberi $n_n^{free}=0$.
        \item Per densit\'a superiori a \SI{1.42e4}{\gram\per\cubic\cm} l'energia di Fermi degli elettroni supera l'energia di legame dell'atomo pi\'u legato $^{56}Fe$ ($Z=26$). Sopra questa densit\'a gli elettroni formano un gas di Fermi libero e i nuclei minimizzano energia formando un lattice.
        \item contribuition of lattice to EOS  is small (see WD): pressure is dominated by degenerate elctron pressure. EOS of outer crust of NS is same of WD.
        \item Energy density is dominated by mass density of nuclei.
        \item Critical density for U R electrons: $E_F=k_F\approx m_e$ quindi $\rho_{rel}=m_un_N=\frac{m_uA}{3\pi^2Z}m_e^3\approx \SI{9.74e5}{\gram\per\cubic\cm}\frac{A}{Z}$.
        \item Change of composition: electron energy arriva a $m_n-m_p$ quindi si ha $\Pproton+Pelectron\to\Pneutron+\Pnue$. $E_F\approx m_n-m_p\approx\SI{1.29}{\mega\ev}$, $\rho_n=\frac{m_uA}{3\pi^2Z}(m_n-m_p)^3=\SI{1.57}{\gram\per\cubic\cm}\frac{A}{Z}$, con densit\'a numerica nucleoni $n_N=\frac{A}{Z}n_p=\frac{A}{Z}n_e$. For $^{56}Fe$ $\rho_n\approx\SI{3.37e7}{\gram\per\cubic\cm}$.
        \item Composizione determinata cercando il nucleo che minimizza energia tot. del sistema in funzione di $(A,Z)$
        \item Il potenziale termodinamico determinante \'e l'energia libera di Gibbs $G(P,T,N)$: \'e particolarmente semplice se consideriamo sistema a (P,T) costanti. Passiamo da U(S,V) con N conservato a G facendo trasformata di Legendre $S\to T$ e $V\to P$: $G=U+PV-TS$ e si ha $dG=VdP-SdT+\mu dN$; usando Eulero $U=-PV+TS+\mu N$ si ha $G=\mu N$: minimizzare G a fisso N coincide con minimizzare $\mu$.
        \item Per $T=0$ e usando le due espressioni ricavate al punto precedente per G si ha $\mu_b=\frac{\epsilon_{tot}+P_{tot}}{n_b}$ con $P_{tot}=n_b^2\PDy{n_b}{\frac{\epsilon_t}{n_b}}|_{(A,Z)}$ e visto che la densit\'a numerica di barioni \'e $n_b=An_i$, infine $P_t=P_e+\frac{1}{3}W_Ln_i$.
        \item EOS can be found minimizing G relto $(A,Z)$
    \end{itemize}
\end{frame}

\begin{frame}{Contribution of lattice to EOS is small: coulomb corrections}
    \begin{itemize}
        \item Spherical Wigner-Seitz cells with radius $r_0$ around nuclei.
        \item Charge neutrality: $Z=\frac{4\pi r_0^3}{3}n_e$.
        \item Coulomb self-energy of charged sphere: $E_{ee}=\int_0^{r_0}dq \frac{q(r)}{r}=\frac{3}{5}\frac{Z^2e^2}{r_0}$, using $q(r)=Ze \frac{r^3}{r_0^3}$
        \item Interaction electrons-ion: $E_i=\int_0^{r_0}dq \frac{Ze}{r}=-\frac{3}{2}\frac{Z^2e^2}{r_0}$ where charge of ion is $Ze$.
        \item $E_{coul}=E_{ee}+E_i=-\frac{9}{10}\frac{Z^2e^2}{r_0}$ hence the coulomb energy per electron from lattice is $\frac{E_{coul}}{Z}=-\frac{9}{10}Z\expy{\frac{2}{3}}e^2(\frac{4\pi}{3}n_e)\expy{\frac{1}{3}}$.
        \item Pressione dovuta al contributo coulombiano: $P_{coul}=n_e^2\TDy{n_e}{(E_{coul}/Z)}|_{(A,Z)}=-\frac{3}{10}Z\expy{\frac{2}{3}}e^2(\frac{4\pi}{3})\expy{\frac{1}{3}}n_e\expy{\frac{4}{3}}$. Politropa con $\Gamma=\frac{4}{3}$.
        \item Coulomb contribution is suppressed compared to Fermi pressure by 1-2 OOM: $Z\expy{\frac{2}{3}}e^2=\alpha Z\expy{\frac{2}{3}}=$ \num{0.01} for He - \num{0.03} for O. ($\alpha=e^2\approx \frac{1}{137}$).
    \end{itemize}
\end{frame}

\begin{frame}{Energia totale della cella di Wigner-Seitz: $\epsilon(A,Z,n_i)=\epsilon_e+n_i(W_L+W_i)$.}
    \begin{itemize}
        \item Considero reticolo bcc  (body centered cubic): ho 1 nucleo al centro della cella e un nucleo in ogni vertice condiviso fra altre 8 cella, quindi in media ho $1+8*\frac{1}{8}=2$ nuclei per ogni cella. Per semplicit\'a ricavo il raggio della sfera di wigner-Seitz attorno a ciascun nucleo: $\frac{2}{a^3}=\frac{1}{\frac{4\pi}{3}r_0}$ quindi $r_0\approx0.5a$.
        \item Neutralit\'a di carica $n_e=Zn_i$
        \item Total energy density of Wigner-Seitz cell: $\epsilon(A,Z,n_i)=\epsilon_e+n_i(W_L+W_i)$, dove $W_L$ is lattice energy and $W_i$ \'e l'energia degli ioni, ovvero la massa. $\epsilon_e$ \'e l'espressione standard per un gas di Fermi $\epsilon_e=\frac{g_e}{(2\pi)^3}\int d^3kE(k)=\frac{g_e}{2\pi^2}\int_0^{K_{F,e}}dkk^2\sqrt{k^2+m^2}$; il contributo del reticolo $W_L$ contiene energia coulombiana degli elettroni e degli ioni: $E_C=E_{ee}+E_i=-f\frac{Z^2e^2}{r_0}$, $f\approx \frac{9}{10}$.
        \item BCC has highest electrostatic potential so is the most stable configuration but difference between other structure is so small that we can have other lattice structure in outer crust of NS.
        \item La struttura del reticolo conta per le propriet\'a di trasporto: conducibilit\'a elettrica, conducibilit\'a termica e scattering elettronico dipendono da disposizione dei nuclei, fononi, vettori del reticolo reciproco, difetti.
        \item lattice poco importante per EO S: $\epsilon_t=\epsilon_i+\epsilon_e+\epsilon_L$, $\epsilon_L\ll\epsilon_i,\epsilon_e$ e anche $P_L=\frac{1}{3}\epsilon_L$ \'e piccola risp $P_e$ quindi i contributi del reticolo influiscono poco su $P(\epsilon)$.
        \item Il contributo de l reticolo \'e importante per determinare quali ioni compongono outer crust: a pressione fissata si minimizza $g(A,Z)=\frac{\epsilon+P}{n_b}$ e il contributo del reticolo \'e determinante per casi $g(A_1,Z_1)\approx g(A_2,Z_2)$.
    \end{itemize}
\end{frame}

\begin{frame}{Most important con tribution total energy of outer crust: nuclear mass $W_N$}
    \begin{itemize}
        \item $W_N=Zm_p+(A-Z)m_n+E_b(A,Z)$. Typical binding energy $-\frac{\SI{8}{\mega\ev}}{A}$; $^{56}Fe$ is the most stable: has lowest total energy per nucleon ($^{62}Ni$ has highest binding energy per nucleon).
        \item Missing nuclei from experimental data, which has to be considered in extreme conditions of NS,  has to be included via nuclear model calculations.
        \item Nuclei having magic number of protons/neutrons are particularly stable: numbers closed shell of protons/neutrons.
    \end{itemize}
\end{frame}

\begin{frame}{Nuclear shell model}
    \begin{figure}[!ht]\includegraphics[trim={0cm 0cm 0 0},rotate=90,clip,keepaspectratio,width=0.95\textwidth]{Schaffner/magic-numbers-shell}\label{fig:Schaffner-magic-numbers-shell}\end{figure}
    Il pattern che emerge dal modello per composizione di Oter Crust of NS \'e analogo a quello dei numeri magici per chiusura shell protoni/neutroni del modello nucleare a shell.
    \begin{columns}[T]
        \begin{column}{0.5\textwidth}
            \begin{itemize}
                \item Posizioni shell chiuse
                \item Dimensione shell gaps
                \item Shell corrections alle masse
            \end{itemize}
        \end{column}
        \begin{column}{0.5\textwidth}
            \begin{itemize}
                \item $S_n$, $S_{2n}$ e drip-line
                \item Q-values delle catture $\beta$
                \item eventualmente livelli e rate di transizione
            \end{itemize}
        \end{column}
    \end{columns}
    
\end{frame}

\begin{frame}{''Modello BPS'' a P fissata in outer crust}
    \begin{itemize}
        \item Per ogni $(A,Z)$ la densit\'a di energia \'e $\epsilon=n_iW_i(A,Z)+\epsilon_e(n_e)+n_iW_L(A,Z,n_i)$ e pressione $P=P_e(n_e)+\frac{1}{3}n_iW_L$, $W_i$ masse nucleari sperimentali o ricavate da modelli nucleari. Fisso $P=P_1$ e cerco il minimo di g al variare di $(A,Z)$.
        \item Prendo nucleo candidato $(A_{11},Z_{11})$ con $n_i=\frac{n_e}{Z_{11}}$ quindi $P=P_e(n_e)+\frac{1}{3}n_iW_L(A_{11},Z_{11},n_e)$: risolvo $P_1=P(A_{11},Z_{11},n_e)$ rispetto a $n_e$ ottenendo $n_e(A_{11},Z_{11},P_1)$ e infine $n_i=\frac{n_e}{Z}$, $n_b=\frac{A}{Z}n_e$.
        \item Calcolo potenziale di Gibbs per barione. A $T=0$ ho $g=\frac{G}{N_b}=\mu_b=\frac{\epsilon+P}{n_b}$ infatti sostituendo $U=TS-PV+\sum_a\mu_aN_a$ in $G=U-TS+PV$ quindi $\sum_a\mu_aN_a=\mu_BN_B+\mu_QN_Q+\mu_LN_L$, inoltre la neutralit\'a impone $N_Q=0$ e per $\nu s$ non intrappolati $\mu_L=0$, infine $g=\frac{G}{n_B}=\frac{V(\epsilon+P)}{n_BV}$ e riscriviamo $g(A,Z;P)=\frac{1}{A}[W_i(A,Z)+Z\mu_e+\frac{4}{3}W_L]$
        \item Determinare minimo di $g$ al variare di $(A,Z)$
        \item Aumentare la pressione: $P_2>P_1$.
        \item Per ogni valore della pressione ho determinato $n_b=\frac{A}{Z}n_e$ e $\rho=m_un_b$.
        \item \xaumenta{P}, \xdiminuisce{\frac{Z}{A}}, compare stato con neutroni liberi energeticamente favorito: esco da modello BPS.
    \end{itemize}
\end{frame}

\begin{frame}{Sequence of nuclei jump discontinuosly: first-order phase transition.}
    \begin{itemize}
        \item Se la sequenza di nuclei passa ad uno pi\'u stabile ho discontinuit\'a nella densit\'a di energia dovuta alla differenza nelle binding energy - first order phase transition. A $T=0$ $\epsilon=n_b\mu_b-P$ e $P, \mu_b$ sono uguali nelle due fasi ma $n_b$ cambia: $\Delta\epsilon=\mu_t\Delta n_b\neq0$.
        \item Costruzione di Maxwell - Al punto di transizione ho $P_{\alpha}=P_{\beta}=P_t$, $\mu_{\alpha}=\mu_{\beta}$: abbiamo transizione da composizione $\alpha$ a $\beta$ a $P$ costante e dato che $P(r)$ nelle stella \'e strettamente monotono si ha una superficie di discontinuit\'a per $(A,Z)$, $n_b$, $\epsilon$.
    \end{itemize}
\end{frame}

\begin{frame}{Modelli nucleari}
    \begin{itemize}
        \item Parametrized mass model: historical originated from liquid drop model (Weizsacker 1935, Bethe and Bacher 1936). Bethe-Weizsacker mass formula
            \begin{equation*}
                B(A,Z)=a_vA+a_{surf}A\expy{\frac{2}{3}}+a_CZ^2A\expy{-\frac{1}{3}}+a_{sym}\frac{(N-Z)^2}{A}+\delta(A)
            \end{equation*}
            where $N=A-Z$, nuclear radius $R=r_0A\expy{\frac{1}{3}}$, $r_0\approx\SIrange{1.0}{1.2}{\femto\meter}$: first term ($\propto R^3$) is the volume term, second term ($\propto R^2$) is the surface term, third term ($\propto \frac{Z^2}{R}$) is proportional to coulomb potential of charges sphere with charge Z e radius R, the fourth term is symmetry term describing isospin dependence of nuclear forces, the last term stands for effect of odd number of protons/neutrons: $\delta(A)=\pm a_{pair}A\expy{-\frac{1}{2}}$ resp for EE and OO numbers of pn and $\delta(A)=0$ if A is odd. Fits the measured nuclear masses (see chowdhury, basu 2006) give coeff: $a_v=-\SI{16}{\mega\ev}$, $a_s=\SI{18}{\mega\ev}$, $a_C=\SI{0.7}{\mega\ev}$, $a_{sym}=\SI{23}{\mega\ev}$, $a_{pair}=\SI{-12}{\mega\ev}$.
        \item Skyrme models: Parametric ansatz for nucleon mean effective potential within \schr{}: potential terms have zero range and can be given in terms of proton/neutron density and gradient thereof. See ''The Skyrme Interaction in finite nuclei and nuclear matter'', ''Skyrme Crystals, Nuclear Matter and Compact Stars'', ''The Skyrme Interaction in finite nuclei and nuclear matter'', ''A unified equation of state of dense matter and neutron star structure'', ''Physics of Neutron Star Crusts''.
        \item Relativistic Mean Field Models - Motivated from meson exchange picture of nuclear forces as proposed by Yukawa and is originated from work of Walecka 1974. Il potenziale del nucleone, nell'equazione di Dirac, \'e mediato dal compo mesonico, a loro volta soluzioni dell'equazione di Klein-Gordon, con sorgente campo mesonico costituito da densit\'a dei nucleoni.
    \end{itemize}
\end{frame}

\begin{frame}{Sequence of nuclei in Outer Crust of NS}
    \begin{columns}[T]
        \begin{column}{0.55\textwidth}
    
    \begin{figure}[!ht]\includegraphics[trim={0cm 0cm 0 0},clip,keepaspectratio,width=0.95\textwidth]{Schaffner/sequence-nuclei-crust-NS}\label{fig:Schaffner/sequence-nuclei-crust-NS}\end{figure}
            
        \end{column}
        \begin{column}{0.45\textwidth}
    \begin{itemize}
        \item Sequence of nuclei in crust of NS: sequence stars at most stable nucleus $^{56}Fe$.
        \item All'aumentare della densit\'a e quindi della densit\'a barionica aumenta anche la densit\'a elettronica per la neutralit\'a di carica (vedi aumento potenziale chimico nelle prime due colonne).
        \item Si raggiunge densit\'a per cui conviene convertire protone e elettrone in neutrone: avremo nuclei neutron rich. Per densit\'a $\rho\approx\SI{8e6}{\gram\per\cubic\cm}$ $^{62}Ni$ \'e energeticamente favorito.
    \end{itemize}
            
        \end{column}
    \end{columns}
            \begin{itemize}
                \item $^{62}Ni$ has magic number $Z=28$; at densities of $\rho\approx\SI{3e9}{\gram\per\cubic\cm}$ appaiono nuclei con numero magico $N=50$.
                \item $^{80}Zn$ is the last nucleus with measured BE. Beyond densities $\rho\ge\SI{6e10}{\gram\per\cubic\cm}$ nuclei are so neutron rich that their masses aren't measures.
                \item Relativistic mean field parameter set TMA predict doubly magic $^{78}Ni$ and that sequence continue along neutron magic number $N=82$
            \end{itemize}
    \begin{itemize}
        \item $^{116}Se$ is predicted to be last nucleus in OC at $\rho\approx\SI{4.55e11}{\gram\per\cubic\cm}$. Here baryon chemical potential have reached neutron mass $\mu_b=m_n=\SI{939.57}{\mega\ev}$: free neutron can be added to the system not bounded to nuclei. ($\mu_b=(\PDy{N_b}{G})_{P,T}$ rappresenta il costo in Gibbs per aggiungere barione al sistema: finch\'e $\mu_b<m_n$ un neutrone libero \'e pi\'u costoso di un barione legato; a densit\'a di neutroni liberi infinitesima un neutrone libero ha $p_{F,n}\to0$ quindi $\mu_n^{free}\to m_nc^2$ e quindi la presenza di un neutrone libero non aumenta Gibbs del sistema, ie $dG_{nuclei}=dG_{free}$).
        \item Neutron drip line marks begins of inner Crust.
    \end{itemize}
\end{frame}

\begin{frame}{Sequence of nuclei in nuclear chart}
    \begin{figure}[!ht]
        \includegraphics[trim={0cm 0cm 0 0},clip,keepaspectratio,width=0.5\textwidth]{Schaffner/nuclearchart-sequence-nuclei-crust-NS}\label{fig:Schaffner-nuclearchart-sequence-nuclei-crust-NS}
    \end{figure}
    \begin{itemize}
        \item Tipicamente in OC abbiamo $Z\approx\numrange{26}{46}$
        \item La carica dei nuclei \'e determinata dal beta-equilibrio: per densit\'a superiori a \SI{6e11}{\gram\per\cubic\cm} si hanno piccoli discostamenti di carica nucleare dovuti alle differenze nei modelli nucleari.
    \end{itemize}<++>
\end{frame}

\begin{frame}{EOS of NS Outer Crust.}
            %lbrt
            \begin{figure}[!ht]\includegraphics[trim={0.4cm 0.4cm 0.3cm 0.5cm},clip,keepaspectratio,width=0.6\textwidth]{Schaffner/EOS-NS-OC}\label{fig:Schaffner-EOS-NS-OC}\end{figure}
    \begin{itemize}
        \item EOS don't shows significative difference for different nuclear models used: nuclear binding energy per nucleone circa \SI{8}{\mega\ev}, massa nucleoni approx \SI{1}{\giga\ev}.
    \end{itemize}
\end{frame}

\begin{frame}{EOS of NS Outer Crust: thickness and mass}
    \begin{align*}
        &\TDy{r}{P(r)}=-\frac{Gm_r(r)\epsilon(r)}{r^2}(1+\frac{P(r)}{\epsilon(r)})(1+\frac{4\pi r^3P(r)}{m_r(r)})\invers{(1-\frac{2Gm_r(r)}{r})}\tag{TOV:HE+GR terms}
    \end{align*}
            %lbrt
    \begin{columns}[T]
        \begin{column}{0.5\textwidth}
            \begin{figure}[!ht]\includegraphics[trim={0.5cm 0.5cm 0.2cm 0.5cm},clip,keepaspectratio,width=0.95\textwidth]{Schaffner/mass-thickness-OC}\label{fig:Schaffner-mass-thickness-OC}\end{figure}
        \end{column}
        \begin{column}{0.5\textwidth}
    \begin{align*}
        &\Rightarrow -\frac{\Delta P}{\Delta R}\approx \frac{P_b}{\Delta R}\approx \frac{GM\epsilon_b}{R^2(1-\frac{2GM}{R})}=\frac{C}{1-2C}\frac{\epsilon_b}{R}    \end{align*}
    \begin{itemize}
        \item As EOS of OC is well determined: solve TOV with modified initial conditions. We assume $\Delta R\ll R$, mass of OC negligible relative to $M_{NS}$, $P_b\ll\epsilon_b$: $\frac{\Delta R}{R}\approx \frac{1-2C}{C}\frac{P_b}{\epsilon_b}$, con compactness $C=\frac{GM}{R}\approx\numrange{0.2}{0.3}$, riscriviamo $\frac{\Delta R}{R}\approx \frac{1-2C}{C}\frac{P_b}{\epsilon_b}$ (GR correction squeeze the crust rel. to newtonian).
    \end{itemize}
        \end{column}
    \end{columns}
    \begin{itemize}
        \item Pressure to energy density ratio. UR \Pelectron: $\mu_e\approx k_{F,e}=(3\pi^2n_e)\expy{\frac{1}{3}}$
            \begin{align*}
                &\frac{P_{e,rel}}{\rho_i}=\frac{(3\pi^2)\expy{\frac{1}{3}}}{4}\frac{n_e\expy{\frac{4}{3}}}{m_uAn_i}=\frac{(3\pi^2)\expy{\frac{1}{3}}}{4}\frac{Zn_e\expy{\frac{1}{3}}}{Am_u}=\frac{Z\mu_e}{Am_u}\Rightarrow \frac{P_b}{\epsilon_b}\approx\num{3e-3}\tag{i.b.: $\frac{Z}{A}=\frac{34}{82}\approx0.3$, $\mu_e\approx\SI{26.29}{\mega\ev}$}\\
                &\Delta R_{OC}=\frac{1-2C}{C}\frac{Z\mu_e}{4Am_u} R\Rightarrow \approx \SI{40}{\meter} (2\msun{}, \SI{10}{\kilo\meter}, C=0.3)\ \SI{200}{\meter} (1.4\msun{}, \SI{15}{\kilo\meter}, C=0.15)\\
                &\Delta M=4\pi\int_0^{\epsilon_b}drr^2\epsilon(r)\approx 4\pi R^2\rho_b\Delta R\approx\num{2.5e-4}\msun{}(\frac{R}{\SI{10}{\kilo\meter}})^2(\frac{\rho_b}{\SI{4e11}{\gram\per\cubic\cm}})(\frac{\Delta R}{\SI{1}{\kilo\meter}})
            \end{align*}
    \end{itemize}
\end{frame}

\subsection{Inner Crust}

\begin{frame}{Inner Crust Phenomenology}
    \begin{itemize}
        \item Inner crust begins at neutron-drip-line, when neutron start being ejected out of nuclei, forming a separate liquid. Dripline is where the baryon chemical potential equals vacuum neutron mass: neutron drip density is model dependent $\rho_{drip}\approx\SIrange{4e11}{5e11}{\gram\per\cubic\cm}$
        \item Relation between baryon density and baryon chemical potential is not unique; baryon chemical potential is energy needed to add/subtract baryon (neutron) from the medium, which depends on potential energy the baryon feels in the medium. For outer crust potential energy of baryon is just binding energy of the nucleus, so the binding energy of last nucleus in the sequence has to be known with precision to determine neutron drip-density.
    \end{itemize}
            \begin{columns}[T]
                \begin{column}{0.45\textwidth}
                    \begin{align*}
                        &g(A,Z;P)=\frac{\epsilon+P}{n_b}\tag{$T=0$}\\
                        &g(A,Z;P)=\frac{1}{A}[W_N(A,Z)+Z\mu_e+\frac{4}{3}W_L]\\
                        &M(A,Z)=Zm_p+(A-Z)m_n-B(A,Z)
                    \end{align*}
                \end{column}
                \begin{column}{0.55\textwidth}
                    \begin{align*}
                        &S_n(A,Z)=M(A-1,Z)c^2+m_nc^2-M(A,Z)c^2\tag{vacuum}\\
                        &S_n<0\tag{emissione neutrone}\\
                        &g(A,Z;P)=\mu_b\tag{OC}
                    \end{align*}
                \end{column}
            \end{columns}
            \begin{itemize}
                \item La struttura della inner crust \'e simile a OC nel fatto che protoni sono all'interno degli ioni organizzati in reticolo circondati da gas di elettroni per bilanciare la carica ma i neutroni sono presenti anche fuori dai nuclei formando un liquido separato.
                \item Il fluido di neutroni \'e superfluido - T. di Cooper (Basis of BCS): degenerate fermions will form bound pairs if there is an attractive potential between them - just one attractive channel of neutron-neutron interaction will make neutron matter superfluid even if the sum of all channels is repulsive. Nucleon scattering experiment: partial wave $^1S_0$ is attractive ($^{2s+1}L_J$), extrapolation of nuclear interaction model to neutron fluid predicts \SI{1}{\mega\ev} attractive potential (at higher density other channels can turn attractive).
                \item Neutron superfluid is corotating (ms-pulsar: rapidly rotating). Rotating superfluid form vortices: quantized angular number momentum trasport.
                    Rotating superfluid is irrotational everywhere except for vortices core. $m\oint v_s\cdot dl=sh$; for uniform rotation $\vec{v}=\vecp{\Omega}{r}$ quindi $\nabla\wedge\vec{v}=\nabla\wedge\vecp{\Omega}{r}=2\vec{\Omega}$, vorticit\'a di un fluido classico in rotazione uniforme.
            \end{itemize}
\end{frame}

\begin{frame}{Vortex structure in inner crust}
    \begin{itemize}
        \item Vorticit\'a di fluido classico in rotazione uniforme $\nabla\wedge\vec{v}=2\vec{\Omega}$; lontano dai vortici $\nabla\wedge\vec{v}_s=0$.
        \item Supponiamo superficie A ortogonale a $\vec{\Omega}$ contenga molti vortici $N_V$ con $s=1$: la circolazione di A \'e la somma di tutti i vortici racchiusi $\oint_{\partial A}\vec{v}_s\cdot d\vec[\mathbf]l=N_V \frac{h}{m}$.
        \item T. Stokes $\oint_{\partial A}\vec{v}_s\cdot d\vec{l}=\int_A\nabla\wedge\vec{v}_s\,dS$: $m\oint_{\partial A}\vec{v}_s\cdot d\vec{l}=N_V \frac{h}{m}=2\Omega A$ quindi $\frac{N_V}{A}=\frac{2\Omega m}{h}$ (Relazione di Feynman per densit\'a dei vortici). Densit\'a momento angolare rispetto asse del vortice $\vec[\mathbf]{l}(\vec{r})=\vec{r}\wedge\rho_s\vec{v}_s$ quindi per vortice assisimmetrico $v_{\phi}(\vec{r})=\frac{s\kappa}{2\pi r}=\frac{s\hbar}{mr}$ dove s \'e il winding number, e $l_z(r)=\rho_srv_{\phi}=\frac{\rho_ss\hbar}{m}$.
        \item I vortici formano reticolo bidimensionale ortogonale all'asse di rotazione: in termini del periodo P della pulsar $\frac{N_V}{A}=\frac{8\pi m_n}{hP}=\SI{6.35e7}{\per\square\meter}\frac{\si{\second}}{P}$. The space between vortices is much smaller than NS size and much bigger than microscopic lattice scale of outer crust.
        \item Pulsar Glitches: Pulsar suddenly rotates faster. Angular momentum transport by vortices. Schematicamente: Spin-down elettromagnetico della crosta, vortici pinned (ostacolati nel moto dalla presenza del reticolo di nuclei), $\Omega_s-\Omega_c$ cresce finch\'e non si ha unpin di molti vortici che si muovono verso l'esterno.
    \end{itemize}
\end{frame}

\begin{frame}{young NS cooling: neutrino emission, meccanismo PBF per $T<T_c$}
    \begin{itemize}
        \item Young NS cool emitting neutrinos. La superfluidit\'a influisce sul raffreddamento di NS attraverso emissione di neutrini. Neutron pairs of bound neutrons have to be broken: per $T<T_c$ in inner crust prevales pairing $^1S_0$ between neutron which produce coherent mixing of many pairs $(k,-k)$ producing BCS condensate. The probability for breaking neutron pair at low $T\neq0$ is suppressed by Boltzmann factor $\exp{-\frac{2\Delta}{kT}}$, which also suppresses neutrino emission rate. Nel superfluido neutronico le eccitazioni elementari sono quasi-particelle di Bogoliubov con $E_k=\sqrt{\xi_k^2+\Delta^2}$, rompere una coppia significa creare due quasi-particelle: condensato $\to\gkupd{k}+\gmkdn{k}$ con gap $2\Delta$. A T finito esiste equilibrio BCS-eccitazioni: $\ketBCS\underset{\text{pair formation}}{\overset{\text{pair breaking}}{\rightleftarrows}}\gkupd{k}\gmkdnd{k}\ketBCS$, $0<T<T_c$. \item Osservazioni della curva di raffreddamento forniscono informazioni su $T_c$ tramite $\Delta(T=0)\approx\num{1.764}kT_c$.
        \item Il \(^{1}S_0\) neutronico della crosta interna \'e  molto importante per la superfluidit\'a e per la dinamica dei vortici, ma nelle curve globali di cooling il gap neutronico del core, soprattutto \(^{3}P_2\!-\!{}^{3}F_2\), \'e generalmente quello che ha generato i vincoli osservativi pi\'u discussi attraverso Cas A e il PBF.
        \item Il pairing ha due effetti fondamentali: sopprime \(C_V\) e molti processi neutrino per \(T\ll T_c\), ma introduce l'emissione PBF quando \(T\) attraversa \(T_c\).
        \item Nel 2011 Page e indipendentemente Shternin interpretarono il rapido calo della temperatura superficiale della giovane NS di Cas A come l'accensione del PBF quando il core neutronico attraversa \(T_c\) per il pairing \(^{3}P_2\). Page et al ottennero una scala \(T_c\sim5\times10^8\,{\rm K}\); Shternin et al. richiedevano un profilo \(T_{cn}(\rho)\) ampio con massimo circa \(7\!-\!9\times10^8\,{\rm K}\). Entrambi richiedevano inoltre protoni gi\'a superconducting, cos\'i da sopprimere parte del cooling precedente.
    \end{itemize}
\end{frame}

\begin{frame}{Cooling delle neutron star e pairing nucleonico}
\small
\begin{itemize}
  \item \textbf{Yakovlev \& Pethick (2004), \emph{Neutron Star Cooling}.}
  Review classica sul raffreddamento delle neutron star:
  processi neutrino, superfluidità, calore specifico e confronto
  fra curve di cooling e osservazioni.
  \item \textbf{Potekhin, Pons \& Page (2015),
  \emph{Neutron Stars---Cooling and Transport}.}
  Review moderna dell'evoluzione termica, trasporto di calore,
  struttura di core/crosta/envelope e interpretazione
  della radiazione termica osservata.
  \item \textbf{Sedrakian \& Clark (2019),
  \emph{Superfluidity in Nuclear Systems and Neutron Stars}.}
  Review microscopica del pairing nucleonico:
  neutroni $^1S_0$ nella crosta, neutroni $^3P_2$--$^3F_2$
  nel core, gap superfluidi, interazioni deboli e vortici.
  \item \textbf{Page et al. (2011),
  \emph{Rapid Cooling of the Neutron Star in Cassiopeia A
  Triggered by Neutron Superfluidity in Dense Matter}.}
  Interpretazione del rapido cooling di Cas A come onset
  del pairing neutronico tripletto e conseguente emissione PBF.
  \item \textbf{Shternin et al. (2011),
  \emph{Cooling neutron star in the Cassiopeia A supernova remnant:
  evidence for superfluidity in the core}.}
  Interpretazione indipendente di Cas A; vincoli sul profilo
  $T_{c,n}(\rho)$ e sul massimo della temperatura critica
  del pairing neutronico nel core.
  \item \textbf{Beloin et al. (2018),
  \emph{Constraining Superfluidity in Dense Matter from the Cooling
  of Isolated Neutron Stars}.}
  Inferenza statistica di $T_c$ neutronico e protonico
  da un insieme di neutron star isolate nel quadro del
  \emph{minimal cooling}.
\end{itemize}
\end{frame}

\begin{frame}{Superfluidit\'a: BCS 101}
    A $T=0$ un gas di fermioni non interagenti riempe tutti gli stati fino $E_F$, $\xi_k=\epsilon_k-\mu$ ($\xi_k<0$ per $\epsilon_k<E_F$): il pairing \'e dominato da stati con $|\xi_k|\ll E_F$. Attrazione debole vicino $E_F$ produce stato legato: Cooper considera due fermioni sopra Fermi's sea inerte - la coppia pi\'u semplice \'e formata da stati time-reversed ($\vec[\mathbf]{k}\uparrow,\vec[\mathbf]{k}\downarrow$), momento totale nullo e pairing in S-wave. Il fermi sea non \'e il GS in presenza di attrazione. (In BCS si considerano si stati sopra, sia stati sotto $E_F$, in cooper solo stati sopra). BCS estende da singola coppia a intero sistema: $\ket{BCS}=\prod_{\mathbf{k}}\left(u_k + v_k \pairkd\right)\ket{0}$, per ogni pair-state k, $u_k^2$ \'e la probabilit\'a che la coppia sia assente, $v_k^2$ la probabilit\'a che sia presente, con $u_k^2+v_k^2=1$. Il parametro di pairing $\kappa_k=\avg{\cmkd}{\cku}$, nel ground-state BCS $\kappa_k=u_kv_k$, \'e l'ampiezza di pairing dello stato di coppia l e attraverso $V_{kl}$ contribuisce al campo efficace sentito da coppia k: il campo di pairing su k $\Delta_k=-\sum_lV_{kl}\kappa_l$. Possiamo ottenere l'espressione precedente considerando $H_{\mathrm{int}}^{{}^1S_0}=\sum_{kl}\underbrace{\left\langle k;{}^1S_0\middle|V_{nn}\middle|l;{}^1S_0\right\rangle}_{V_{kl}}\,\underbrace{b_k^\dagger b_l}_{\text{pair scattering}}$, dove $\bpair{l}=\cmkdn{l}\ckup{l}$. Introduco mean-field approx $\bpair{k}\approx\overbrace{\avg{\bpair{k}}}^{\kappa_k}+\delta\bpair{k}$, dinamica collettiva complicata→ogni grado di libert\'a vede il campo medio degli altri, con $\delta\bpaird{k}=\bpaird{k}-\avg*{\bpaird{k}}$, $\delta\bpair{l}=\bpair{l}-\avg*{\bpair{l}}$ quindi $\bpaird{k}\bpair{l}=\avg*{\bpaird{k}}\avg*{\bpair{l}}+\delta\bpaird{k}\,\avg*{\bpair{l}}+\avg*{\bpaird{k}}\,\delta\bpair{l}+\delta\bpaird{k}\,\delta\bpair{l}$ quindi $H_{int}^{MF}=-\sum_k(\Delta_k\bpaird{k}+\Delta_k^*\bpair{k})+\sum_k\kappa_k^*\Delta_k$, a $H_{int}^{MF}$ aggiungiamo $\sum_{k\sigma}\xi_k\,c^\dagger_{k\sigma}c_{k\sigma}=H_0-\mu\hat N$ con $\xi_k=\epsilon_k-\mu$ e nel modello di particelle libere $\epsilon_k=\frac{\hbar^2k^2}{2m}$, termine a una particella, per avere $H^{MF}$: il primo termine risulta da $\sum_{kl}V_{kl}\bpaird{k}\kappa_l=\sum_k\bpaird{k}(\sum_lV_{kl}\kappa_l)=-\sum_k\Delta_k\bpaird{k}$; per il secondo termine considero interazione hermitiana $V_{kl}=V_{lk}^*$, quindi $\Delta_k^*=-\sum_lV_{lk}^*\kappa_k^*=-\sum_lV_{kl}\kappa_k^*$, nel modello BCS semplice $\Delta_k=\Delta$. Nel normale gas di Fermi l'energia rispetto a $\mu$ \'e $|\xi_k|$ mentre ne lsuperfluido BCS $E_k=\sqrt{\Delta^2+\xi_k^2}$; alla superficie di Fermi $\xi_k=0$ quindi $E_{min}=\Delta$. Le eccitazioni sono combinazioni particle/hole $\gamma_k=u_kc_k-v_kc_{-k}^{\dag}$, quasi-particelle di Bogoliubov: si ha che $u_k^2=\frac{1}{2}(1+\frac{\xi_k}{E_k})$, $v_k^2=\frac{1}{2}(1-\frac{\xi_k}{E_k})$ con $u_kv_k=\frac{\Delta}{2E_k}$, vicino $k_F$ $u_k^2\approx v_k^2\approx \frac{1}{2}$, quindi il mixing particle-hole \'e massimo. Il pairing smussa l'occupazione attorno a $E_F$, quindi tende ad aumentare l'energia cinetica, ma abbassa maggiormente l'energia di interazione. Il risultato netto \'e $U_s-U_n = -\frac{1}{2}N(0)\Delta^2$ nel weak-coupling BCS semplice.
\end{frame}

\begin{frame}{Inner Crust EOS}
    \begin{itemize}
        \item Reticolo di celle di Wigner–Seitz, ciascuna contenente un cluster nucleare ricco di neutroni immerso in un gas di neutroni ''drippati'' e in un fondo elettronico: Negele e Vautherin (1973) trattano ogni cella come una sfera di raggio \(R_{\rm WS}\) e risolvono in modo self-consistent Hartree–Fock il problema nucleare nella cella: cercano le densit\'a $n_p(r)$, $n_n(r)$ che minimizzano l'energia per data densit\'a media della materia, variando anche composizione $(N,Z)$ e dimensione della cella.
        \item I protoni hanno picco di densit\'a al centro della cella, i neutroni hanno picco al centro della cella e profilo piatto non nullo verso l'esterno; all'aumentare della densit\'a aumenta la densit\'a neutronica nella parte di piatta, infine si restringe anche la cella e il profilo neutronico della densit\'a \'e costante con piccolo incremento centrale.
        \item Strong shell effects for protons: for lower densities number of protons inside the cell is $Z=40$, at higher densities switches to $Z=50$: a magic number for the shell model without a strong spin-orbit splitting and to a closed subshell for the case of strong spin-orbit splitting; proton numbers staying around \numrange{40}{50} is a balance between the surface energy and the Coulomb energy of the Wigner-Seitz cell; the total neutron number inside a cell reaches $N=1800$.
        \item Inner crust stops at $0.5n_0\approx\SI{0.075}{\per\cubic\femto\meter}$ ($0.5\epsilon_0=0.5m_nn_0\approx\SI{70}{\mega\ev\per\cubic\femto\meter}$), lattice size shrinks from \SI{100}{\femto\meter} to \SI{30}{\femto\meter}.
        \item For higher densities, denpending upon nuclear model, we reach outer core consisting of liquid neutrons with small admixture of protons, electrons; lattice disappears.
    \end{itemize}
\end{frame}

\begin{frame}{Hartree--Fock autoconsistente nella cella di Wigner--Seitz}

\small
\begin{itemize}
  \item Si approssima una cella della crosta interna con una sfera di raggio
  \(R_{\mathrm{WS}}\), assumendo simmetria sferica:
  \[
    n_n(\mathbf r)\to n_n(r),
    \qquad
    n_p(\mathbf r)\to n_p(r).
  \]

  \item Nel metodo Hartree--Fock il ground state è approssimato da un determinante
  di Slater e gli orbitali soddisfano
  \[
    h_{\mathrm{HF}}[n_n,n_p]\phi_i=\epsilon_i\phi_i .
  \]

  \item L'autoconsistenza consiste nel richiedere che il campo medio costruito
  dalle densità produca, risolvendo le equazioni HF, le stesse densità:
  \[
    n_q(r)=\sum_{i\in q}^{\mathrm{occ}}|\phi_i(r)|^2,
    \qquad q=n,p .
  \]

  \item Schema iterativo:
  \[
  n_q^{(0)}(r)
  \rightarrow
  h_{\mathrm{HF}}[n_q^{(0)}]
  \rightarrow
  \phi_i
  \rightarrow
  n_q^{(1)}(r)
  \rightarrow
  \cdots
  \rightarrow
  n_q^{(m+1)}(r)\simeq n_q^{(m)}(r).
  \]

  \item La configurazione della cella si ottiene minimizzando l'energia totale
  rispetto a
  \[
    N,\qquad Z,\qquad R_{\mathrm{WS}} .
  \]
\end{itemize}

\end{frame}

\begin{frame}{Stima spessore crosta e pasta}
    \begin{itemize}
        \item For inner pressure we take pressure of NR neutrons, instead of pressure of UR electron as in outer crust. At inner boundary of inner crust (gas degenere NR, $T=0$)
            \begin{columns}[T]
                \begin{column}{0.42\textwidth}
            \begin{align*}
                &\frac{P_b}{\epsilon_b}|_{crust}\approx\xi \frac{P_n^{NR}}{\rho_n}=\frac{\xi}{5}(3\pi^2)\expy{\frac{2}{3}}\frac{n_n\expy{\frac{2}{3}}}{m_n^2}=\frac{\xi}{5}\frac{k_{F,n}^2}{m_n^2}\\
                &E_{F,n}=\frac{k_{F,n}^2}{2m_n}\Rightarrow \frac{P_b}{\epsilon_b}\approx \frac{2\xi}{5}\frac{E_{F,n}}{m_n}\\
                &\approx \frac{0.4}{5}(\frac{257}{939})^2\approx\num{6e-3}
            \end{align*}
                \end{column}
                \begin{column}{0.58\textwidth}
                    $P_b=\xi P_n^{NR}$, $\xi=0.4$, al fondo della crosta $\epsilon_b\approx\rho_n\approx m_nn_n$, $k_{F,n}=(3\pi^2n_n)\expy{\frac{1}{3}}$, densit\'a saturazione nucleare $n_0=\SI{0.15}{\per\cubic\femto\meter}$, ma nella materia nucleare $n_n\approx n_p\approx \frac{n_0}{2}$. La pressione \'e piccola rispetto alla densit\'a di energia perch\'e $E_F\ll m_m$. Per $n_n=n_0$ si ha $k_{F,n}\approx\SI{1.64}{\per\femto\meter}$, usando $\hbar c\approx\SI{197}{\mega\ev\femto\meter}$ si ha $k_{F,n}\approx\SI{324}{\mega\ev}$; a $\frac{n_0}{2}$ $K_{F,n}\approx\SI{257}{\mega\ev}$
                \end{column}
            \end{columns}
                   
        \item \xaumenta{M_*}, \xdiminuisce{\Delta R_{crust}} - Usando la relazione $\frac{\Delta R}{R}\approx \frac{1-2C}{C}\frac{P_b}{\epsilon_b}$: \SI{80}{\meter} for a $2\msun{}$ star with $R=\SI{10}{\kilo\meter}$ and $C=0.3$, and \SI{400}{\meter} for a $1.4\msun{}$ star with $R=\SI{15}{\kilo\meter}$ and $C=0.15$.
        \item Pasta in the Crust (Ravehall '83)- Asymmetrical Wigner-Seitz cells: cells gettin closer can marge in one direction to for a rod of nuclear matter, when this happen for several chains along same direction we have a spaghetti-like structure surrounded by liquid neutron. Increasing further density many rods form a plate, giving rise to a lasagna-like structure of nuclear matter plates alternating with pure neutron matter plate. Aumentando ancora la densit\'a si invertono i ruoli tra nuclear matter densa e dripping di liquido neutronico: la lasagna \'e self-inverted, per gli spaghetti avro rods di neutroni liquidi circondati da materia nucleare densa, infine bolle di neutroni liquidi circondati da materia nucleare densa.
        \item Condizioni stabilit\'a pasta-phase: considero energia singola cella
            \begin{columns}[T]
                \begin{column}{0.25\textwidth}
                    \begin{align*}
                        &\frac{E_{tot}}{V}=\epsilon_t=\epsilon_b+\epsilon_s+\epsilon_C+\epsilon_e
                    \end{align*}
                \end{column}
                \begin{column}{0.75\textwidth}
                    Densit\'a energia cella: bulk, superficie, Coulomb, elettroni. Energia minimizzata rispetto al raggio $R_N$ della struttura nucleare: bulk e electron non dipendono da $R_N$ quindi non influenzano il minimo: $n_e(r)$ uniforme (lunghezza onda compton $\frac{1}{2\pi m_e}\approx\SI{e-12}{\meter}$, $R_{WS}\approx\SIrange{10}{20}{\femto\meter}$)
                \end{column}
            \end{columns}
            
    \end{itemize}
\end{frame}
